\section{Criação do Conjunto de Dados}
\label{sec:criacao_dados}

Ocorpus abrange diferentes subdomínios da Cibersegurança, incluindo segurança da informação e sistemas de Inteligência Artificial. O conjunto final é constituído por 22 pares pergunta-resposta, dos quais 12 foram elaborados manualmente e 10 gerados automaticamente. Cada registo inclui o tema, a pergunta, o método de geração (HUM ou LLM) e a respetiva resposta validada

\subsection{ Geração manual}
Foram criadas manualmente 12 perguntas sobre diferentes tópicos da Cibersegurança.

\subsection{Geração Automática}
\label{sec:geracao_automatica}
A geração automática foi realizada com o modelo AMALIA q4\_K\_M, disponibilizado através da plataforma Ollama \cite{amalia_ollama}.
O prompt foi executado seis vezes  para aumentar a diversidade das perguntas geradas.

\textit{prompt} utilizado

\begin{quote}
\small
Gera uma lista de perguntas factuais em português sobre o tema Cibersegurança,
combinando conceitos clássicos de segurança da informação (criptografia,
protocolos de rede, controlo de acessos) e segurança em sistemas de Inteligência
Artificial (vulnerabilidades em LLMs, ataques a modelos).

Cada pergunta deve ser direta, clara, autossuficiente e ter uma resposta curta e
inequívoca, formulada no máximo numa frase. Evita perguntas ambíguas, demasiado
abertas ou que exijam explicações longas.

Não repitas perguntas nem o mesmo conceito de forma redundante. Procura abranger diferentes áreas da cibersegurança.

Apresenta cada resultado exatamente no seguinte formato:

Pergunta: [pergunta]

Resposta: [resposta curta]
\end{quote}


